Research on emergent misalignment reports many null results, and each is read as evidence that an intervention did not install a misaligned pathway. None of those nulls is equivalence-tested, so none separates ``the pathway is absent'' from ``this measurement did not elicit it.'' We propose \mspfull (\msp): the minimum number of training or evaluation decisions that must be flipped simultaneously, away from a null-producing baseline specification, before misalignment exceeds a pre-registered threshold $\tau$. We measure \msp over a full binary specification lattice on \qwenseven with a local \qwenfourteen judge, splitting it into an elicitation arm (four evaluation-side factors, $2^4$ cells) and an induction arm (three training-side factors). Our headline result is negative and, we argue, consequential: the elicitation-side \msp is degenerate at $1$ for all three conditions we test, \emph{including the unmodified base model}, where a single persona flip moves misalignment from $0.0\%$ ($0/22$) to $93.8\%$ ($15/16$, $p=5.6\times 10^{-14}$). The same flip takes a 5\%-dilution ``null'' organism to $100\%$ ($8/8$). Once each cell must also exceed the base model at the identical evaluation specification, no flip-set in the lattice qualifies for either finetuned condition. Separately, finetune-induced activation shifts are low-dimensional (effective rank $5.8$--$7.6$ against a random-subspace null of $14.9$), which argues for a bottleneck rather than a redundancy account. We conclude that eval-side elicitation cannot diagnose a misalignment null without a base-model control, and we give the first equivalence-tested nulls in this literature.
